\chapter{Criterio de demo seleccionada}

\section{Demo seleccionada}
\begin{definition}
Se llama demo seleccionada a toda pieza metodológicamente admitida que, dentro del Tomo III, es elegida para exposición, seguimiento y eventual traducción disciplinada a claim por su claridad estructural, su trazabilidad y su valor ejemplificador.
\end{definition}

\begin{observation}
Una demo seleccionada no equivale todavía a resultado final cerrado. Su función es servir como pieza metodológicamente tratable, suficientemente clara y útil para evaluación, exposición o proyección disciplinada.
\end{observation}

\section{Condición mínima de selección}
\begin{definition}
Una pieza puede ser tratada como demo seleccionada cuando concurren simultáneamente:
\begin{enumerate}
  \item admisión metodológica ya satisfecha;
  \item función ejemplificadora o curatorial reconocible;
  \item trazabilidad suficiente de origen, nivel y transformaciones;
  \item forma de presentación compatible con seguimiento disciplinado;
  \item posibilidad de traducción posterior a un uso demostrativo o selectivo más fuerte.
\end{enumerate}
\end{definition}

\begin{observation}
La selección no exige todavía cierre final de validación. Exige solo que la pieza ya sea útil como soporte metodológico claro y no como material bruto o prematuro.
\end{observation}

\section{Demo y configuración}
\begin{definition}
Una configuración operativa puede convertirse en demo seleccionada cuando su función, su dependencia y su recorrido pueden presentarse como caso tratable y legible dentro del Tomo III.
\end{definition}

\begin{observation}
Esta cláusula evita seleccionar configuraciones cuya identidad siga siendo demasiado opaca para uso metodológico.
\end{observation}

\section{Demo y familia}
\begin{definition}
Una familia operativa puede convertirse en demo seleccionada cuando sus miembros permiten mostrar un patrón común disciplinado sin destruir las diferencias internas relevantes para el propósito declarado.
\end{definition}

\begin{observation}
La selección de una familia no exige uniformidad total. Exige solo cohesión suficiente para exposición y seguimiento.
\end{observation}

\section{Demo y traducción}
\begin{definition}
Una traducción entre niveles puede convertirse en demo seleccionada cuando explicita de forma legible qué conserva, qué modifica y por qué sigue siendo metodológicamente utilizable.
\end{definition}

\begin{observation}
Esta cláusula impide seleccionar traducciones cuyo valor dependa de supuestos implícitos no declarados.
\end{observation}

\section{Condición mínima de no selección}
\begin{definition}
Se considera que una pieza no debe ser seleccionada como demo cuando, aun siendo admisible, carece de claridad estructural suficiente, de trazabilidad legible o de valor metodológico reconocible para el propósito declarado.
\end{definition}

\begin{observation}
Esta condición protege al Tomo III frente a selección inflada, arbitraria o meramente acumulativa.
\end{observation}

\section{Cláusula de alcance}
\begin{observation}
El criterio de demo seleccionada fija una primera regla de elección metodológica dentro del Tomo III. No constituye todavía una teoría completa de evidencia final ni una selección definitiva de claims.
\end{observation}

